% =====================================================================
\section{Mock Exam -- Form A}\label{ch:exam-a}
% =====================================================================

\textbf{Instructions.} Answer all problems. Show complete solutions; box your
final answers. Simplify exact answers (no decimals unless asked).
Time: 2 hours. Total: 100 points.

\subsection*{Part I. Vector functions (35 points)}

\begin{questions}
\item Find the domain of
      $\vr(t)=\ip{\sqrt{4-t^2},\;\ln(t+1),\;\dfrac{1}{t-1}}$. \pts{5}
\item Let $C$ be the curve of intersection of the cylinder $x^2+y^2=4$ and the
      plane $x+z=3$.
      \begin{enumerate}[label=(\alph*)]
        \item Find a vector function that represents $C$.
        \item Find parametric equations of the tangent line to $C$ at $(0,2,3)$.
      \end{enumerate}\pts{10}
\item Find $\vr(t)$ if $\vr\,'(t)=\ip{2t,\;3t^2,\;e^t}$ and $\vr(0)=\ip{1,0,-1}$. \pts{5}
\item Find the length of the curve $\vr(t)=\ip{2t,\;t^2,\;\ln t}$, $1\le t\le e$. \pts{7}
\item Find the curvature of the twisted cubic $\vr(t)=\ip{t,t^2,t^3}$ at the point $(1,1,1)$. \pts{8}
\end{questions}

\subsection*{Part II. Directional derivatives and the gradient (35 points)}

\begin{questions}[resume]
\item Find the directional derivative of $f(x,y)=x^3-2xy^2$ at $P(1,-1)$ in the
      direction toward $Q(4,3)$. \pts{8}
\item Let $T(x,y,z)=x^2+y^2-2z^2$ and $P(1,1,1)$.
      \begin{enumerate}[label=(\alph*)]
        \item Find the rate of change of $T$ at $P$ in the direction toward the origin.
        \item In which direction does $T$ increase most rapidly at $P$? What is the maximum rate?
        \item Find an equation of the tangent plane to the level surface of $T$ through $P$.
      \end{enumerate}\pts{15}
\item Let $f(x,y)=x^2+xy$ and $P(1,1)$. Find all unit vectors $\vu$ such that
      $\Du f(P)=1$. Is there a direction in which $\Du f(P)=4$? Explain. \pts{12}
\end{questions}

\subsection*{Part III. Vector calculus (30 points)}

\begin{questions}[resume]
\item Sketch the vector field $\vF(x,y)=\vi+x\,\vj$ by drawing its vectors at
      the nine points $(x,y)$ with $x,y\in\{-1,0,1\}$. Describe the field. \pts{6}
\item Determine whether $\vF(x,y)=\ip{2x+y\cos(xy),\;x\cos(xy)+3y^2}$ is
      conservative. If it is, find a potential function. \pts{8}
\item Show that $\vF(x,y,z)=\ip{2xz+y^2,\;2xy,\;x^2+3z^2}$ is conservative and find
      a potential function $f$ with $f(0,0,0)=5$. \pts{9}
\item Find the value of the constant $a$ for which
      $\vF=\ip{axy-z^3,\;(a-2)x^2,\;(1-a)xz^2}$ is conservative, and for that value
      find a potential function. \pts{7}
\end{questions}

\begin{center}\textit{--- End of Form A ---}\end{center}

% =====================================================================
\subsection{Solutions to Form A}

\solhead{1}{5}
Each component must be defined.
\begin{itemize}
\item $\sqrt{4-t^2}$: $4-t^2\ge0\iff-2\le t\le2$.
\item $\ln(t+1)$: $t+1>0\iff t>-1$.
\item $\dfrac1{t-1}$: $t\neq1$.
\end{itemize}
Intersection: $\boxed{(-1,1)\cup(1,2]}$.\quad\textit{[2 pts for the three conditions, 3 for the intersection.]}

\solhead{2}{10}
\textbf{(a)} The cylinder gives $x=2\cos t$, $y=2\sin t$. From the plane,
$z=3-x=3-2\cos t$. Hence
\[
  \vr(t)=\ip{2\cos t,\;2\sin t,\;3-2\cos t},\qquad0\le t\le2\pi .
\]
Check: $x^2+y^2=4$\checkok, $x+z=3$\checkok.

\textbf{(b)} The point $(0,2,3)$: $2\cos t=0$, $2\sin t=2\Rightarrow t=\frac\pi2$, and
$z=3-0=3$\checkok. $\vr\,'(t)=\ip{-2\sin t,\;2\cos t,\;2\sin t}$, so
$\vr\,'\!\left(\frac\pi2\right)=\ip{-2,0,2}$. Tangent line:
\[
  \boxed{x=-2s,\quad y=2,\quad z=3+2s.}
\]

\solhead{3}{5}
Integrate componentwise: $\vr(t)=\ip{t^2,\;t^3,\;e^t}+\vb C$. At $t=0$:
$\ip{0,0,1}+\vb C=\ip{1,0,-1}$, so $\vb C=\ip{1,0,-2}$ and
$\boxed{\vr(t)=\ip{t^2+1,\;t^3,\;e^t-2}}$.

\solhead{4}{7}
$\vr\,'(t)=\ip{2,\;2t,\;\frac1t}$, so
\[
  \norm{\vr\,'(t)}=\sqrt{4+4t^2+\frac1{t^2}}=\sqrt{\left(2t+\frac1t\right)^2}=2t+\frac1t\qquad(t>0),
\]
since $\left(2t+\frac1t\right)^2=4t^2+4+\frac1{t^2}$. Then
\[
  L=\int_1^e\left(2t+\frac1t\right)\dd t=\Big[t^2+\ln t\Big]_1^e=(e^2+1)-(1+0)=\boxed{e^2}.
\]
\textit{[Recognizing the perfect square is the key step: 3 pts.]}

\solhead{5}{8}
$(1,1,1)$ corresponds to $t=1$. $\vr\,'=\ip{1,2t,3t^2}$, $\vr\,''=\ip{0,2,6t}$; at $t=1$:
$\vr\,'=\ip{1,2,3}$, $\vr\,''=\ip{0,2,6}$.
\[
  \vr\,'\times\vr\,''=\begin{vmatrix}\vi&\vj&\vk\\1&2&3\\0&2&6\end{vmatrix}
  =(12-6)\vi-(6-0)\vj+(2-0)\vk=\ip{6,-6,2},
\]
$\norm{\vr\,'\times\vr\,''}=\sqrt{76}=2\sqrt{19}$ and $\norm{\vr\,'}^3=(\sqrt{14})^3=14\sqrt{14}$. So
\[
  \kappa(1)=\frac{2\sqrt{19}}{14\sqrt{14}}=\frac{\sqrt{19}}{7\sqrt{14}}=\boxed{\frac{\sqrt{266}}{98}}\approx0.166 .
\]

\solhead{6}{8}
$f_x=3x^2-2y^2$, $f_y=-4xy$ (continuous, so $f$ is differentiable);
$\grad f(1,-1)=\ip{3-2,\;4}=\ip{1,4}$. $\vecPQ=\ip{3,4}$, $\norm{\vecPQ}=5$,
$\vu=\ip{\frac35,\frac45}$. Hence
$\Du f(1,-1)=\frac35+\frac{16}{5}=\boxed{\frac{19}{5}}$.

\solhead{7}{15}
$\grad T=\ip{2x,2y,-4z}$, so $\grad T(1,1,1)=\ip{2,2,-4}$.

\textbf{(a)} Toward the origin: $\ip{-1,-1,-1}$, $\vu=\frac{1}{\sqrt3}\ip{-1,-1,-1}$.
$\Du T(P)=\frac{-2-2+4}{\sqrt3}=\boxed{0}$. \textit{[The direction is tangent to the
level surface; see (c): the origin lies on the cone $x^2+y^2=2z^2$.]}

\textbf{(b)} Direction of $\ip{2,2,-4}$, i.e. unit vector $\frac{1}{\sqrt6}\ip{1,1,-2}$;
maximum rate $\norm{\ip{2,2,-4}}=\sqrt{24}=\boxed{2\sqrt6}$.

\textbf{(c)} $T(P)=1+1-2=0$: the level surface is the cone $x^2+y^2-2z^2=0$, with
normal $\grad T(P)$:
$2(x-1)+2(y-1)-4(z-1)=0\iff\boxed{x+y-2z=0}$.

\solhead{8}{12}
$\grad f=\ip{2x+y,\;x}$, so $\grad f(1,1)=\ip{3,1}$, $\norm{\grad f}=\sqrt{10}$.
For $\vu=\ip{a,b}$: $3a+b=1$ and $a^2+b^2=1$. Then $b=1-3a$ and
$a^2+(1-3a)^2=1\iff10a^2-6a=0\iff a=0$ or $a=\frac35$.
\[
  \boxed{\vu=\ip{0,1}\quad\text{or}\quad\vu=\ip{\tfrac35,-\tfrac45}}
\]
(check: $0+1=1$, $\frac95-\frac45=1$\checkok). $\Du f(P)=4$ is impossible because
$\abs{\Du f(P)}\le\norm{\grad f(P)}=\sqrt{10}\approx3.16<4$.

\solhead{9}{6}
$\vF(x,y)=\ip{1,x}$ depends only on $x$:
\[
  x=-1:\ \ip{1,-1},\qquad x=0:\ \ip{1,0},\qquad x=1:\ \ip{1,1}
\]
at each of the three heights $y=-1,0,1$.
\begin{center}
\begin{tikzpicture}[scale=1.1,>=Stealth]
  \draw[->,gray] (-2,0) -- (2.6,0) node[right]{$x$};
  \draw[->,gray] (0,-1.8) -- (0,1.9) node[above]{$y$};
  \foreach \x in {-1,0,1}{ \foreach \y in {-1,0,1}{
     \fill[NavyBlue] (\x,\y) circle (1.5pt);
     \draw[->,thick,BrickRed] (\x,\y) -- ({\x+0.5},{\y+0.5*\x});
  }}
  \node[below left] at (0,0) {\small $0$};
\end{tikzpicture}
\end{center}
(arrows drawn at half length). Every arrow points to the right; arrows tilt
down for $x<0$, are horizontal on the $y$-axis and tilt up for $x>0$, the tilt
growing with $\abs{x}$; all arrows on a vertical line are identical.

\solhead{10}{8}
$M=2x+y\cos(xy)$, $N=x\cos(xy)+3y^2$.
\[
  M_y=\cos(xy)-xy\sin(xy),\qquad N_x=\cos(xy)-xy\sin(xy).
\]
Equal and continuous on $\R^2$: \textbf{conservative}. Then
$f=\int M\,\dd x=x^2+\sin(xy)+g(y)$ (since $\pd{}{x}\sin(xy)=y\cos(xy)$);
$f_y=x\cos(xy)+g'(y)=x\cos(xy)+3y^2\Rightarrow g=y^3+K$.
\[
  \boxed{f(x,y)=x^2+\sin(xy)+y^3+K}.
\]

\solhead{11}{9}
\textbf{Test:} $M_y=2y=N_x$;\ $N_z=0=P_y$;\ $P_x=2x=M_z$. Conservative.\\
\step{1}{$x$.} $f=x^2z+xy^2+g(y,z)$.
\step{2}{$y$.} $f_y=2xy+g_y=2xy\Rightarrow g=h(z)$.
\step{3}{$z$.} $f_z=x^2+h'(z)=x^2+3z^2\Rightarrow h=z^3+K$.\\
$f(0,0,0)=K=5$, so $\boxed{f=x^2z+xy^2+z^3+5}$.

\solhead{12}{7}
$M_y=ax$ and $N_x=2(a-2)x$; equal iff $a=2a-4$, i.e. $a=4$.
$N_z=0=P_y$ always. $P_x=(1-a)z^2$ and $M_z=-3z^2$: equal iff $a=4$.\checkok\
So $\boxed{a=4}$ and $\vF=\ip{4xy-z^3,\;2x^2,\;-3xz^2}$.
$f=2x^2y-xz^3+g(y,z)$; $f_y=2x^2\Rightarrow g_y=0$; $f_z=-3xz^2+h'(z)=-3xz^2\Rightarrow h=K$.
$\boxed{f=2x^2y-xz^3+K}$.
